\section{Введение}

\begin{frame}[c]{Зачем нужно машинное обучение?}
  \lead{Вопрос на подумать.}
  Алан Тьюринг расшифровывал Энигму перебором вариантов --- машина справилась быстрее людей.

  \bigskip
  \begin{block}{Есть задачи, которые человеку даются легко, а программой их почти не записать}
    Назовите примеры таких задач.
  \end{block}
\end{frame}

\begin{frame}[c]{Ответ: легко человеку, трудно запрограммировать}
  \begin{columns}[T]
    \begin{column}{0.52\textwidth}
      \begin{numrows}
        1 & Перевести текст с языка на язык\\
        2 & Поставить диагноз по симптомам\\
        3 & Выбрать, какой из двух документов лучше подходит запросу\\
        4 & Сказать, что изображено на картинке\\
        5 & Оценить, по какой цене продастся квартира\\
      \end{numrows}\par
    \end{column}
    \begin{column}{0.44\textwidth}
      \vspace*{-4pt}
      \begin{block}{Что у них общего}
        \begin{enumerate}\setlength\itemsep{0.3em}
          \item Решение --- функция из объекта в предсказание.
          \item Единственно верного ответа нет: переводы бывают разными, врачи ошибаются.
          \item Правильных ответов много, а неправильные легко придумать.
        \end{enumerate}
      \end{block}
    \end{column}
  \end{columns}
\end{frame}

\begin{frame}[c]{Что такое машинное обучение?}
  \vfill
  \lead{Вопрос.}

  \medskip
  {\Large\bfseries\color{acDark} Объясните своими словами, как человеку, далёкому от IT.\par}
  \vfill
\end{frame}

\begin{frame}[c]{Определение машинного обучения}
  \lead{Компьютер сам находит закономерности в данных.}
  Вместо того чтобы человек вручную прописывал правила, алгоритм выучивает зависимость между
  объектами и ответами по примерам и применяет её к новым объектам.

  \medskip
  \begin{columns}[T]
    \begin{column}{0.31\textwidth}
      \begin{block}{Данные}
        Примеры: объекты и правильные ответы к ним.
      \end{block}
    \end{column}
    \begin{column}{0.31\textwidth}
      \begin{block}{Модель}
        Функция $f(x) \to y$: по объекту $x$ выдаёт предсказание $y$.
      \end{block}
    \end{column}
    \begin{column}{0.31\textwidth}
      \begin{block}{Обучение}
        Подбор параметров модели так, чтобы ответы на данных были как можно лучше.
      \end{block}
    \end{column}
  \end{columns}
\end{frame}

\begin{frame}[c]{Из чего состоит задача}
  \centering
  \begin{tikzpicture}[
      box/.style={rectangle,minimum width=0.27\textwidth,minimum height=0.9cm,
        text=white,font=\bfseries,align=center},
      arr/.style={-{Stealth[length=2.5mm]},line width=1.2pt,color=acGray}]
    \node[box,fill=acDark] (x) {Объект $x$};
    \node[box,fill=acMid,right=0.8cm of x] (f) {Модель $f$};
    \node[box,fill=acDark,right=0.8cm of f] (y) {Предсказание $y$};
    \draw[arr] (x) -- (f);
    \draw[arr] (f) -- (y);
  \end{tikzpicture}

  \medskip
  \begin{columns}[T]
    \begin{column}{0.31\textwidth}
      \begin{block}{Объекты и таргеты}
        Объект --- картинка, история болезни, пользователь сервиса.
        Таргет --- ответ: подпись, диагноз, уход пользователя.
      \end{block}
    \end{column}
    \begin{column}{0.31\textwidth}
      \begin{block}{Модель}
        Функция $f(x) \to y$. Принимает признаки объекта $x$ и выдаёт предсказание $y$.
      \end{block}
    \end{column}
    \begin{column}{0.31\textwidth}
      \begin{block}{Датасет}
        Обучающая выборка: объекты вместе с известными ответами. По ней учится модель.
      \end{block}
    \end{column}
  \end{columns}
\end{frame}
